%% JASA submission, generated by make_jasa.py from ../main.tex.
\documentclass[12pt]{article}

\usepackage[margin=1in]{geometry}
\usepackage{amsmath,amssymb,amsthm}
\usepackage{graphicx}
\usepackage{booktabs}
\usepackage{adjustbox}
\usepackage{enumerate}
\usepackage[round,authoryear]{natbib}
\usepackage{xurl}
\usepackage[hidelinks]{hyperref}

% 12pt leaves less room for inline formulas: allow looser lines, not overfull ones
\setlength{\emergencystretch}{3em}

\newcommand{\E}{\mathbb{E}}
\newcommand{\KL}{\mathrm{KL}}
\newcommand{\F}{\mathcal{F}}
\newcommand{\G}{\mathcal{G}}

\newtheorem{theorem}{Theorem}
\newtheorem{proposition}{Proposition}
\newtheorem{corollary}{Corollary}
\theoremstyle{remark}
\newtheorem{remark}{Remark}

% 0: version with the author; 1: anonymised version for double-blind review.
\newcommand{\blind}{0}

\def\spacingset#1{\renewcommand{\baselinestretch}%
{#1}\small\normalsize} \spacingset{1}

\begin{document}

\if0\blind
{
  \title{\bf Betting against the simulator:\\
  Monte Carlo e-processes under adaptive actions}
  % Funding, if any, goes in a \thanks{...} after the name.
  \author{Elena Bortolato\\
    CUNEF Universidad, Madrid, Spain}
  \date{}
  \maketitle
} \fi

\if1\blind
{
  \bigskip
  \bigskip
  \bigskip
  \begin{center}
    {\LARGE\bf Betting against the simulator:\\[4pt]
    Monte Carlo e-processes under adaptive actions}
  \end{center}
  \medskip
} \fi

\bigskip
\begin{abstract}
Multi-arm multi-stage trials, bandits and digital twins collect data
adaptively. We test whether a simulator reproduces such a system with
e-processes, which allow optional stopping. A score learned by classifying real
against simulated outcomes becomes an exact bet when normalised over $M$ fresh
simulations. Under adaptive allocation, validity follows from exchangeability
if the simulator is run under the action taken and from the current state.
This conditioning is necessary: a simulator that ignores the action makes the
expected bet tend to $1+\chi^2$, and the wealth grows exponentially under the
null. The actions also add power: in an audit of a language-model API,
choosing the prompts adaptively saves every score up to 61\% of the queries.
The growth rate of any score splits into a simulation cost, $M/(M+1)$ for local
alternatives, and a learning cost, $\KL(Q\|\tilde T)$, with $\tilde T$ the
score's forecast. The bet is a likelihood in the score's parameter, so Bayesian
updating is exact: posterior-predictive bets stay pathwise within $d\log t$ of
the best score in hindsight. In a multi-arm multi-stage trial of insulin
dosing, the monitor controls false alarms and identifies the failing arms;
monitors that ignore the allocation or the participants' history do not.
\end{abstract}

\noindent%
{\it Keywords:} e-values; anytime-valid inference; simulation-based inference;
digital twins; adaptive designs
\vfill

\newpage
\spacingset{1.45}
\section{Introduction}\label{sec:intro}
%======================================================================

A digital twin is useful only as long as it behaves like the system it
copies. Checking this is hard when the system is being steered while we
watch. Consider a multi-arm multi-stage (MAMS) trial of insulin-dosing
strategies \citep{whitehead1997design,magirr2012generalized}. At each interim
analysis some arms are dropped and their participants move to other arms,
taking their glucose history with them. A simulator of each participant's
glucose dynamics is available, and the question is whether, and when, the
real participants stop behaving like simulated ones.

The same question comes up with multi-armed bandits
\citep{lattimore2020bandit}, response-adaptive clinical designs
\citep{trippa2012bayesian,villar2015multi}, and AI systems that act in
environments shaped by their own past decisions
\citep{perdomo2020performative,dulacarnold2021challenges}. For example, a ranking or moderation system changes what users post and hence
what the platform observes next \citep{chaney2018algorithmic,jiang2019degenerate},
and such systems are tuned on simulators of user behaviour whose fidelity is
exactly what is in doubt \citep{ie2019recsim,shi2019virtual,zhao2020sim}. In these settings, as we choose
actions based on past observations, can we continuously test whether the
resulting stream of outcomes is still consistent with a simulator, and sound
an alarm as soon as it is not, while controlling false alarms under optional
stopping? The data are dependent, because what is observed next depends on
what has been seen; the null is a simulator whose likelihood we
cannot write down; and we want to look at the evidence after
every observation. Fixed-sample tests lose
their guarantees under repeated looks, and the corrections developed for
adaptive trials \citep{bowden2017unbiased,magirr2012generalized} assume a
pre-specified adaptation rule. Bayesian updating copes with adaptive sampling
but does not, by itself, control the type~I error under optional stopping.

E-values suit this situation
\citep{vovk2021evalues,ramdas2023game,grunwald2024safe}. An e-process is the
running wealth of a gambler who places one fair bet against the null per
observation. By Ville's inequality, the probability that this wealth ever
exceeds $1/\alpha$ is at most $\alpha$, so the experimenter can look after
every observation and stop at any time. The best known e-process is the Bayes
factor. Written one observation at a time, it is a product of ratios: the
posterior predictive density of each new observation over its null density.
Each factor has mean one under the null for every prior, and the product
rewards whoever forecasts the data better, the learner or the null
\citep{dawid1984present}. Learning the parameters along the way costs only a
term of order $\log t$ \citep{clarke1990information}.

We cannot use the Bayes factor here, for two reasons. The null density is not
available, since the simulator has latent dynamics and, under adaptive
allocation, its law changes with the path. Nor is there a trusted model of the
real system to forecast with, since that is precisely what is in question.
Our approach keeps the shape of the Bayes factor, one fair bet per
observation placed before the observation is seen, and replaces both of its
ingredients.

To replace the null density, we start from a classifier. Trained to
distinguish real outcomes from simulated ones, a classifier estimates the
ratio of their densities from samples alone \citep{thomas2022lfire}. Call
\emph{score} any positive function $T$ of an outcome that is large where real
outcomes are more likely than simulated ones; the classifier's estimate of the
density ratio is the natural choice. A score is not a fair bet by itself, but
a Monte Carlo test turns it into one \citep{besag1989generalized}. At each
step $t$ we observe the real outcome $x_t$, run the simulator $M$ times to
obtain $y_{t,1},\dots,y_{t,M}$, and bet $M+1$ times the real outcome's share
of the total score. If the simulator is right, the real outcome is just one
more draw among $M+1$ exchangeable ones, and the bet is fair for any score. In
an adaptive design the only delicate point is how to simulate. We run the
simulator under the action actually taken and from the current state of the
unit. This is enough for validity under any allocation rule, and it cannot be
dropped: a simulator run without regard to the action produces false alarms
with probability one.

The score also replaces the forecast. Multiplying the simulator's density by
the score and normalising gives a density for the next outcome, which moves
the simulator's mass towards where the score expects real outcomes to fall:
choosing a score amounts to forecasting reality. The score is unknown and must
be learned as the data arrive, so we let it depend on parameters $\theta$, for
example $T_\theta(x) = \exp\{\theta^\top z(x)\}$ for a vector of features
$z(x)$. The bet is itself a likelihood for $\theta$. Up to the factor $M+1$,
the bet is the real outcome's share of the total score,
$T_\theta(x_t)/\{T_\theta(x_t) + \sum_{j} T_\theta(y_{t,j})\}$: the
probability that a conditional-logit model \citep{mcfadden1974conditional},
which picks one item from a set
with probability proportional to its score, picks the real outcome among the
$M+1$ candidates. The product of these probabilities over time is a likelihood
for $\theta$, so $\theta$ can be learned by Bayes' rule: we put a prior on
$\theta$, update the posterior after each observation, and place the next bet
with the posterior average of the bet. Since the posterior uses only past
observations, this average is still a fair bet under the null. It plays the
role of the posterior predictive in the Bayes factor, with the score in place
of the model. Standard regret bounds for Bayesian mixtures
\citep{kakade2005online} then guarantee that, on every sample path, the wealth
stays within a $d\log t$ term of what the best score in the class would have
earned in hindsight.

Our contributions are the following.
\begin{enumerate}
\item \emph{Validity under adaptive actions, and why the conditioning is
needed.} The Monte Carlo e-process is a test martingale for every allocation
rule (Theorem~\ref{thm:valid}). Without conditioning on the action, the
expected bet under the null tends to $1+\chi^2$ for the likelihood-ratio score
and false alarms grow exponentially, at a rate of at most $\log K$ per step
with $K$ arms (Proposition~\ref{prop:necessary},
Corollary~\ref{cor:falsealarm}).
\item \emph{Bayesian updating of the score.} Posterior-predictive bets,
updated in one step per observation, with a regret bound on every sample path
(Theorem~\ref{thm:simlik}); approximate posteriors keep validity exact at a
bounded cost in power (Proposition~\ref{prop:approx}).
\item \emph{Adaptive actions as a source of power.} Since validity holds for
any allocation rule, the actions can be chosen to collect evidence quickly. In
an audit of a language-model API, choosing the prompts adaptively reduces the
number of queries by up to 61\%, for our scores and for those of concurrent
methods (Section~\ref{sec:exp-llm}).
\item \emph{The power of an arbitrary score.} A score implicitly bets that
reality looks like the simulator reweighted by the score, a tilted version of
the null. Its growth rate is the Kullback--Leibler information against the
simulator minus the distance from reality to this tilted null, and $M$
simulations keep a fraction $M/(M+1)$ of it for local alternatives
(Theorem~\ref{thm:growth}).
\end{enumerate}

Validity comes from how the simulator is run, not from the score, so any score
fixed before the new outcome is seen, including a network trained in advance,
gives a valid test. Our update rules are meant for the usual situation in which
nothing is known about the alternative: a twin tends to fail in ways its
designers did not anticipate.

In the insulin trial, with five strategies over three stages, only the monitor
that runs the twin from each participant's state under the assigned arm
controls false alarms (2\% of trials, against 95\% for a twin that ignores the
participants' history). One e-process per arm, combined by e-BH
\citep{xu2021unified}, identifies the arm on which the twin fails.

Three recent papers use the same Monte Carlo bet for independent data
\citep{dombowsky2026besag,forre2026anytime,cortinovis2026slice}; for a single
test it is the soft-rank e-value of \citet{wang2022false}. We developed our
construction independently (Section~S9 of the Supplement). Their main
developments assume independent observations and a simulator without state.
\citet{forre2026anytime} list dependent data as an open problem, and
\citet[Appendix~C.1.1]{cortinovis2026slice} note that the construction extends
under conditional exchangeability and prove validity given independent
covariates. They also learn the score differently: \citet{forre2026anytime}
refit it by maximising its empirical growth, and \citet{cortinovis2026slice}
train a neural ratio estimator before the data arrive, over the simulator's own
parametric family, so that their score can forecast only departures that the
family contains; our scores reweight the simulator by features of the outcome,
and detect departures that its designers did not build into it.
Section~\ref{sec:discussion} lists in detail which of our results coincide
with theirs.

Section~\ref{sec:method} gives the method and its validity,
Section~\ref{sec:power} its power, and Section~\ref{sec:predictive} the
predictive view and the update rules (Algorithm~S1 in the Supplement). Section~\ref{sec:experiments}
contains simulations, the comparison with concurrent methods and the
language-model audit, and Section~\ref{sec:trial} the trial.
Section~\ref{sec:discussion} discusses limitations.

%======================================================================
\section{Monte Carlo e-processes under adaptive control}\label{sec:method}
%======================================================================

\subsection{Setting}\label{sec:setting}

The symbols used throughout are collected in Table~S11 of the
Supplement.

At each time $t = 1, 2, \ldots$ an action (or context) $A_t$ is chosen,
possibly at random, by a policy that may depend on everything observed
before, including the monitor's own wealth. The real system then produces an
outcome $X_t$ with conditional law $Q_t(\cdot \mid \F_{t-1}, A_t)$, and the
analyst can draw samples from a simulator with conditional law
$P_t(\cdot \mid \F_{t-1}, A_t)$; for a digital twin this means re-initialising
the twin at the current state and running it under action $A_t$. The
filtration $\F_t$ contains the actions, the real outcomes, the simulator draws
and any internal randomisation up to time $t$, and
$\G_t = \F_{t-1} \vee \sigma(A_t)$ is the information available when the bet
at time $t$ is placed. The null hypothesis is
\[
H_0:\quad Q_t(\cdot \mid \F_{t-1}, A_t) = P_t(\cdot \mid \F_{t-1}, A_t)
\quad \text{for all } t, \text{ almost surely}:
\]
the simulator reproduces the conditional behaviour of the system under every
action that is actually taken. $H_0$ is composite in the policy, which is left
unrestricted. Because each factor is fair conditionally on the past and on
the action, one e-process is valid for every policy at once, with no infimum
over the class. The policy may depend
on past outcomes and on the monitor's own wealth; the price is that
the null concerns only the actions actually taken, so the monitor cannot
detect failures of the simulator under actions the policy never chooses. A nonnegative process $(W_t)$ with $W_0 = 1$ is a \emph{test
supermartingale} for $H_0$ if $\E[W_t \mid \F_{t-1}] \le W_{t-1}$ under every
distribution in $H_0$; by Ville's inequality, rejecting the first time $W_t$
reaches $1/\alpha$ controls the type-I error at level $\alpha$ uniformly over
time and stopping rules \citep{ramdas2023game,grunwald2024safe}. Products of
nonnegative factors with
\begin{equation}\label{eq:fair}
\E[E_t \mid \G_t] \le 1 \quad \text{under } H_0
\end{equation}
are test supermartingales.

\subsection{The Monte Carlo e-process}

Let $s_t$ be a real-valued score chosen from the information in $\G_t$, and
$T_t = e^{s_t}$. Draw $Y_{t,1}, \ldots, Y_{t,M}$ from the simulator under
action $A_t$ and set
\begin{equation}\label{eq:mc}
E_t = \frac{(M+1)\, T_t(X_t)}{T_t(X_t) + \sum_{j=1}^{M} T_t(Y_{t,j})}.
\end{equation}

\begin{proposition}[One-step exactness; \citealp{wang2022false}; \citealp{dombowsky2026besag}, Proposition~1]\label{prop:exact}
Let $U_0, \ldots, U_M$ be exchangeable and $T > 0$. Then
$\E\big[(M+1)\,T(U_0)/\sum_{k} T(U_k)\big] = 1$.
\end{proposition}

\begin{proof}
The ratios $T(U_k)/\sum_{k'} T(U_{k'})$ sum to one and, by exchangeability,
have equal expectations, each equal to $1/(M+1)$.
\end{proof}

The normalised scores $T_t(U_k)/\sum_{j} T_t(U_j)$ can be viewed as the
probabilities that the
score assigns to each of the $M+1$ outcomes being the real one, and $E_t$ is
$M+1$ times the probability it assigns to the outcome that is actually real.
The bet pays when the score picks out the real outcome better than a uniform
guess, $1/(M+1)$; under the null the real outcome is indistinguishable from
the simulations, so on average no score can.

A score is typically a classifier's estimate of the ratio $q_t/p_t$ of the
densities of the real outcome and of the simulator (Section~\ref{sec:scores}).
Used directly as the bet, $T_t(X_t)$ (the \emph{plug-in} bet) has conditional
mean $Z_t = \E_{P_t}[T_t]$ under the null, which equals one only if the
estimated ratio integrates exactly to one against the simulator. A learned
ratio does not. If $\log T_t = \log(q_t/p_t) + \varepsilon$ with
$\E[\varepsilon] = 0$, then $\E[Z_t] = \E[e^{\varepsilon}] > 1$ by Jensen's
inequality even when $q_t = p_t$, so estimation noise alone makes the plug-in
bets favourable under the null; this is why an overfitted learner gives many
false alarms (Section~\ref{sec:exp-drift}). The factor \eqref{eq:mc} replaces
$Z_t$ by its estimate
$\widehat Z_t = \{T_t(X_t) + \sum_{j=1}^{M} T_t(Y_{t,j})\}/(M+1)$, and
Proposition~\ref{prop:exact} shows that this self-normalisation makes the
conditional mean exactly one, whatever $Z_t$ is.

The factor is bounded, $0 < E_t \le M+1$. To hedge the fraction of wealth that
is bet, we average the processes $\prod_i (1 - \lambda + \lambda E_i)$ over a
finite grid $\Lambda \subset (0, 1]$,
\begin{equation}\label{eq:wealth}
W_t = \frac{1}{|\Lambda|} \sum_{\lambda \in \Lambda} \prod_{i \le t}
\big(1 - \lambda + \lambda E_i\big).
\end{equation}

\begin{theorem}[Anytime validity under adaptive control]\label{thm:valid}
Suppose that, conditionally on $\G_t$, the simulator draws
$Y_{t,1}, \ldots, Y_{t,M}$ are i.i.d.\ from $P_t(\cdot \mid \F_{t-1}, A_t)$ and
independent of $X_t$, and that the score $s_t$ is $\G_t$-measurable. Then
under $H_0$, $\E[E_t \mid \G_t] = 1$ for every $t$, the wealth
\eqref{eq:wealth} is a test martingale, and
$P(\exists t: W_t \ge 1/\alpha) \le \alpha$. This holds for every policy,
every learner and every $M \ge 1$.
\end{theorem}

\begin{proof}
Under $H_0$ and conditionally on $\G_t$, the real outcome $X_t$ has law
$P_t(\cdot \mid \F_{t-1}, A_t)$ and is independent of the simulator draws, so
$(X_t, Y_{t,1}, \ldots, Y_{t,M})$ is i.i.d.\ and hence exchangeable, while
$T_t$ is fixed. Proposition~\ref{prop:exact} gives
$\E[E_t \mid \G_t] = 1$, so $\E[1 - \lambda + \lambda E_t \mid \F_{t-1}] = 1$
for each $\lambda$. Each product is a nonnegative martingale with initial
value one, and so is their average. Ville's inequality concludes.
\end{proof}

For an i.i.d.\ null without actions, Theorem~\ref{thm:valid} reduces to the
Besag--Clifford e-process of \citet[Proposition~6]{dombowsky2026besag}. The
theorem involves no density, no property of the score and no restriction on
the policy, which may depend on the data and on the wealth itself; in a MAMS
trial the action $A_t$ is the arm assigned to a participant for a period,
chosen after interim analyses and re-assignments, and the conditional law
$P_t$ is that of the participant's outcome given their history. The policy
enters only through the requirement that the simulator be run under the action
actually taken. Proposition~\ref{prop:necessary} shows that this
requirement cannot be dropped.

\begin{proposition}[Conditioning on the action cannot be dropped]\label{prop:necessary}
Suppose that $H_0$ holds and that at time $t$ the policy takes an action $a$,
with $P_a = P_t(\cdot \mid \F_{t-1}, a)$, while, conditionally on $\G_t$, the
simulator draws $Y_{t,1}, \ldots, Y_{t,M}$ are i.i.d.\ from a law $\bar P$
that does not depend on $A_t$ (for instance the twin run under uniformly
random actions) and are independent of $X_t$. Let $T > 0$ be a bounded
$\G_t$-measurable score.
\begin{enumerate}
\item[(i)] \emph{The limit of the expected bet.} As $M \to \infty$,
\[
\E[E_t \mid \G_t] \to \rho = \frac{\E_{P_a}[T]}{\E_{\bar P}[T]}
\quad\text{and}\quad
\E[E_t^2 \mid \G_t] \to \frac{\E_{P_a}[T^2]}{(\E_{\bar P}[T])^2}.
\]
\item[(ii)] \emph{Favourable scores exist whenever $P_a \ne \bar P$.} The
limit $\rho$ equals one for every bounded score if and only if $P_a = \bar P$.
If $P_a \ll \bar P$ and $R_a = dP_a/d\bar P$, the likelihood-ratio score
gives
\[
\E_{P_a}[R_a] = \int R_a^2\, d\bar P = 1 + \chi^2(P_a \,\|\, \bar P):
\]
$T = R_a$ has $\rho = 1 + \chi^2(P_a\|\bar P)$ when $R_a$ is bounded and
bounded away from zero, and in general the truncated scores
$T_c = \min(R_a, c) + 1/c$ have $\rho \to 1 + \chi^2(P_a\|\bar P) \in (1,\infty]$
as $c \to \infty$.
\item[(iii)] \emph{Hence false alarms.} Suppose that $P_a$ and $\bar P$ do not
change over time and that, on an event $\Omega_0$, from some finite time on
the policy plays $a$ with a fixed score $T$ taking values in $[b, B]$,
$0 < b \le B$, whose limit in (i) is $\rho > 1$; by (ii) such a score exists
whenever $P_a \ne \bar P$. Then there is $M_0$ such that, for every
$M \ge M_0$, if $\Lambda$ contains some
$\lambda \le \min\{1/2,\, (\rho - 1)b^2/(8B^2)\}$, on $\Omega_0$, almost
surely,
\[
\liminf_{t\to\infty} \frac{\log W_t}{t} \ge \frac{\lambda(\rho - 1)}{4} > 0 .
\]
In particular the test rejects the true null with probability at least
$P(\Omega_0)$.
\end{enumerate}
\end{proposition}

The proof is in Section~S7 of the Supplement. The event $\Omega_0$ occurs, for
example, when a bandit settles on the best arm: each bet then has mean larger
than one, by (i) and (ii), and the wealth grows exponentially, by (iii). The
proposition also shows that the condition $P_a = \bar P$ is not only
sufficient for fairness (Theorem~\ref{thm:valid}) but necessary for the bet to
be fair against every score, and that $1 + \chi^2(P_a\|\bar P)$ is the exact
limit of the expected bet of the likelihood-ratio score.

For a single Monte Carlo test, simulations drawn from a law within total
variation $\varepsilon$ of the right one inflate the level by at most
$\varepsilon$ \citep[Theorem~2]{barber2026montecarlo}. In a sequential
e-process the bias of every factor compounds instead: the expected bet is
governed by the $\chi^2$ divergence and the growth of the wealth by the
Kullback--Leibler divergence (Corollary~\ref{cor:falsealarm}), and
$2\,\mathrm{TV}^2 \le \KL \le \log(1+\chi^2)$, so any departure, however small in
total variation, eventually raises a false alarm. Moreover $\chi^2$ can be
large when the total variation is small, if the real system produces rare
outcomes that the simulator almost never does.

%======================================================================
\section{Power: the growth rate}\label{sec:power}
%======================================================================

The time needed to reach $1/\alpha$ is governed
by the growth rate $\E[\log E_t \mid \G_t]$ \citep{grunwald2024safe}. Fix a
time, condition on $\G_t$ and drop the index. Let $P$ be the simulator
law and $Q$ the true law of the next outcome, with $Q \ll P$; let $T = e^{s}$
and let $Z = \E_P[T] < \infty$ be the normalising constant of
Section~\ref{sec:method}. The score reweights the simulator law into a new
law for the next outcome, its \emph{tilted null},
\begin{equation}\label{eq:tilt}
\frac{d\tilde T}{dP} = \frac{T}{Z},
\end{equation}
which keeps the outcomes of the simulator and gives each a weight
proportional to its score: a change of measure on the space of outcomes, not
the distribution of the score values. It has density $Tp/Z$ when $P$ has
density $p$, and we use the same symbol for the law and its density. It is the
score's forecast of $Q$, and a perfect score has $\tilde T = Q$
(Section~\ref{sec:scores} develops this reading). Write
$\rho = \E_Q[T]/Z$, the limit of the expected bet as $M \to \infty$
(Proposition~\ref{prop:necessary}(i), with $Q$ and $P$ in place of $P_a$
and $\bar P$). Let $X \sim Q$
and, independently, $Y_1, \ldots, Y_M$ i.i.d.\ $P$, and let
$E^{(M)} = (M+1)T(X)/\{T(X) + \sum_{j} T(Y_j)\}$ be the factor \eqref{eq:mc}
computed from them. The \emph{growth rate} with $M$ simulator draws is
\begin{equation}\label{eq:growthrate}
\begin{aligned}
g_M = g_M(Q, P, T) &= \E_{Q \otimes P^M}\big[\log E^{(M)}\big] \\
&= \int \log\frac{(M+1)\,T(x)}{T(x) + \sum_{j=1}^M T(y_j)}\;
Q(dx) \prod_{j=1}^M P(dy_j),
\end{aligned}
\end{equation}
the expected logarithm of one factor when the whole wealth is bet. Since
$\log E^{(M)} \le \log(M+1)$ it is well defined in $[-\infty, \log(M+1)]$, and
finite when $s$ is bounded. It is at most zero when $Q = P$, by Jensen's
inequality and Proposition~\ref{prop:exact}. When a fraction $\lambda$ of the
wealth is bet, as in each term of \eqref{eq:wealth}, the rate is
\[
g_M(\lambda) = g_M(\lambda; Q, P, T)
= \E_{Q \otimes P^M}\big[\log(1 - \lambda + \lambda E^{(M)})\big],
\qquad g_M = g_M(1).
\]
In the sequential procedure $Q$, $P$ and $T$ are the conditional laws and the
score at time $t$, and the \emph{conditional growth rate}
$g_{M,t} = \E[\log E_t \mid \G_t]$ is $g_M$ evaluated at them; it is a
random variable that depends on the history. We write $g_M$ without arguments
when $Q$, $P$ and $T$ are clear from the context.

\begin{theorem}[Growth rate]\label{thm:growth}
Assume $\KL(Q\|P) < \infty$ and $\rho < \infty$.
\begin{enumerate}
  \item[(i)] (Ceiling.) $g_M \le \KL(Q\|P)$ for every $M$ and every score.
  \item[(ii)] (Unlimited budget.) As $M \to \infty$,
        $E^{(M)} \to (d\tilde T/dP)(X)$ almost surely, and if $s$ is bounded,
        $g_M \to g_\infty = \E_Q[\log (d\tilde T/dP)(X)] = \KL(Q\|P) - \KL(Q\|\tilde T)$.
  \item[(iii)] (Finite budget.) $g_M \ge g_\infty - \log\{1 + (\rho - 1)/(M+1)\}
        \ge g_\infty - (\rho - 1)/(M+1)$.
  \item[(iv)] (Local alternatives.) If $dQ/dP = 1 + \varepsilon \psi$ and
        $T = e^{\varepsilon h}$ with $\psi, h$ bounded and $\E_P[\psi] = 0$, and
        $\tilde h = h - \E_P[h]$, then
        $g_M = \frac{M}{M+1} \cdot \frac{\varepsilon^2}{2}
        \big(\E_P[\psi^2] - \E_P[(\psi - \tilde h)^2]\big) + O(\varepsilon^3)$,
        while $\KL(Q\|P) = \frac{\varepsilon^2}{2}\E_P[\psi^2] + O(\varepsilon^3)$.
\end{enumerate}
\end{theorem}

Part (i) is the log-optimality of the likelihood ratio. When the score is the
likelihood ratio, $\tilde T = Q$ and (ii) is the limit in Theorem~1 of
\citet{dombowsky2026besag}; part (iii) is their Theorem~3, rewritten with the
tilted null. What is new is the reading of an arbitrary score through
$\tilde T$ in (ii), and the local constant in (iv), which factorises the
efficiency relative to the oracle likelihood ratio into a simulation term and
a learning term,
\[
\frac{g_M}{\KL(Q\|P)} \approx
\underbrace{\frac{M}{M+1}}_{\text{simulation}} \times
\underbrace{\Big(1 - \frac{\KL(Q\|\tilde T)}{\KL(Q\|P)}\Big)}_{\text{learning}}.
\]
One draw per observation halves the growth rate and so roughly doubles the time
to detection; nine draws keep 90\% and sixteen 94\%. Away from the local regime,
$M/(M+1)$ is neither an upper nor a lower bound: for a large mean shift and a
perfect score the efficiency at $M = 1$ is lower, while an overconfident score
does better with a finite budget than with an infinite one, because
$E^{(M)} \le M+1$ truncates its bets (Section~\ref{sec:exp-budget}).

Theorem~\ref{thm:growth} describes one factor. Over the whole sequence, the
log-wealth is at least the sum of the conditional growth rates of its factors,
up to fluctuations that average out, so the test rejects once that sum exceeds
$\log(1/\alpha)$.

\begin{corollary}[Cumulative growth]\label{cor:cumulative}
For $\lambda \in \Lambda$ let
$g_{M,i}(\lambda) = \E[\log(1 - \lambda + \lambda E_i) \mid \G_i]$ be the
conditional growth rate of the $\lambda$-term of \eqref{eq:wealth} at time $i$.
\begin{enumerate}
\item[(i)] For every $\lambda \in \Lambda$ with $\lambda < 1$, and with no
condition on the scores,
$\log W_t \ge \sum_{i \le t} \log(1 - \lambda + \lambda E_i) - \log|\Lambda|$,
each increment lies in $[\log(1-\lambda),\, \log(1 + \lambda M)]$, and almost
surely
\[
\liminf_{t \to \infty} \frac1t \Big(\log W_t - \sum_{i \le t} g_{M,i}(\lambda)\Big) \ge 0 .
\]
\item[(ii)] If $1 \in \Lambda$ and the scores are bounded, $|s_t| \le B$, the
same holds for $\lambda = 1$, with increments $\log E_i \in [-2B, \log(M+1)]$
and rates $g_{M,i} = g_{M,i}(1)$.
\end{enumerate}
If, for some $\lambda \in \Lambda$, the average of $g_{M,i}(\lambda)$ over
$i \le t$ stays above $\bar g > 0$, the test rejects with probability one,
after about $(\log(1/\alpha) + \log|\Lambda|)/\bar g$ steps.
\end{corollary}

\begin{proof}
The average \eqref{eq:wealth} is at least $|\Lambda|^{-1}$ times any one of its
terms. Since $0 < E_i \le M+1$ by \eqref{eq:mc}, the increment
$\log(1 - \lambda + \lambda E_i)$ lies in $[\log(1-\lambda), \log(1+\lambda M)]$
for $\lambda < 1$, whatever the score; for $\lambda = 1$ and $|s_i| \le B$,
every term of \eqref{eq:mc} lies between $e^{-B}$ and $e^{B}$, so
$\log E_i \in [-2B, \log(M+1)]$. In both cases the differences
$\log(1 - \lambda + \lambda E_i) - g_{M,i}(\lambda)$ are bounded martingale
differences for the filtration $(\G_{i+1})$, and the martingale strong law
(Azuma--Hoeffding and Borel--Cantelli, as in the proof of
Proposition~\ref{prop:necessary}) shows that their average tends to zero
almost surely.
\end{proof}

 Since $1-\lambda+\lambda E$ is a weighted average of 1
and $E$, concavity of the logarithm gives $\log(1-\lambda+\lambda E) \ge \lambda\log E$,
so $g_{M,i}(\lambda) \ge \lambda\, g_{M,i}$: betting a fraction $\lambda$ keeps
at least a fraction $\lambda$ of the growth rate. Part (ii) needs
$|s_t| \le B$, which gives $\log E_i \ge -2B$; clipping the score keeps
validity (Theorem~\ref{thm:valid}), and our learners clip at $B = 30$. A ridge
penalty alone does not bound the score: the estimate can grow like
$\sqrt t$, and polynomial features are unbounded.

The same theory gives the rate of false alarms of a twin that ignores the
action. Theorem~\ref{thm:growth} depends only on the laws of the real outcome
and of the simulations, which for such a twin are $P_a$ and $\bar P$ even
under the true null; the monitor therefore grows as if it were testing $P_a$
against $\bar P$.

\begin{corollary}[False-alarm rate of an action-ignoring twin]\label{cor:falsealarm}
In the setting of Proposition~\ref{prop:necessary}, suppose that the actions
take finitely many values, that the laws $P_a$ and $\bar P$ and the score
$T_a = e^{s_a}$ used when $a$ is played, with $|s_a| \le B$, do not change
over time, that $\KL(P_a\|\bar P) < \infty$ for every action $a$, and that the
fraction of the first $t$ steps at which the policy plays $a$ converges almost
surely to $\varpi_a$. Let $g_M(a) = g_M(P_a, \bar P, T_a)$ be the growth rate
\eqref{eq:growthrate} with the real law $P_a$ in the role of $Q$ and the
simulator law $\bar P$ in the role of $P$, $g_\infty(a)$ its limit as
$M \to \infty$, and
$\tilde T_a = T_a\bar P/\E_{\bar P}[T_a]$. Then, under $H_0$:
\begin{enumerate}
\item[(i)] $g_\infty(a) = \KL(P_a\|\bar P) - \KL(P_a\|\tilde T_a)$ and
$g_M(a) \ge g_\infty(a) - (\rho_a - 1)/(M+1)$, where
$\rho_a = \E_{P_a}[T_a]/\E_{\bar P}[T_a]$ is the limit of
Proposition~\ref{prop:necessary}(i). If $1 \in \Lambda$, almost surely
$\liminf_t t^{-1}\log W_t \ge \sum_a \varpi_a\, g_M(a)$.
\item[(ii)] For every $\lambda$, every $M$ and every score,
$g_M(\lambda; P_a, \bar P, T_a) \le \KL(P_a\|\bar P)
\le \log\{1 + \chi^2(P_a\|\bar P)\}$, and almost surely
$\limsup_t t^{-1}\log W_t \le \sum_a \varpi_a \KL(P_a\|\bar P)$. The
likelihood-ratio score attains the bound as $M \to \infty$.
\item[(iii)] If the twin is run under randomised actions,
$\bar P = \sum_b u_b P_b$ with $u_b > 0$, then
$\KL(P_a\|\bar P) \le \log(1/u_a)$ and $\chi^2(P_a\|\bar P) \le 1/u_a - 1$;
with $K$ actions drawn uniformly, $\log K$ and $K - 1$.
\end{enumerate}
\end{corollary}

The proof is in Section~S7 of the Supplement. The false alarm therefore arrives
after roughly $\log(1/\alpha)/\sum_a \varpi_a g_M(a)$ steps: the rate is a
Kullback--Leibler divergence, not the $\chi^2$ divergence that governs the
expected bet, and it is largest for a policy that concentrates on one action,
$\varpi_a \to 1$, as Thompson sampling does. The rate vanishes only if the arms
have the same law, and with $K$ arms it can never exceed $\log K$ per step.

%======================================================================
\section{Betting is predicting}\label{sec:predictive}
%======================================================================

By Theorem~\ref{thm:growth}, a good score is one whose tilted null is close
to the truth. Choosing such a score is therefore a problem of
sequential prediction: before each real outcome is seen, forecast it as well
as possible from the past. When the null density is known, this problem has a
classical solution, the Bayes factor, and we build on it. We first recall why
the Bayes factor is a product of forecasts and show that the Monte Carlo
e-process has the same structure, with the score in place of the model
(Section~\ref{sec:scores}); we then derive an update rule for the score with a
pathwise guarantee (Section~\ref{sec:simlik}), and its practical version
(Section~\ref{sec:blr}).

For a model $q_\theta$ with prior $\pi_0$, the Bayes factor is the ratio of
the marginal likelihood $m(x_{1:t}) = \int \prod_{i \le t} q_\theta(x_i \mid x_{1:i-1})\,
\pi_0(d\theta)$ to the null density $p_0$. By the chain rule,
\begin{equation}\label{eq:prequential}
\mathrm{BF}_t = \prod_{i=1}^{t} \frac{m_i(x_i)}{p_0(x_i \mid x_{1:i-1})},
\qquad
m_i(x) = \int q_\theta(x \mid x_{1:i-1})\,\pi_{i-1}(d\theta),
\end{equation}
where $\pi_{i-1}$ is the posterior after $i-1$ observations, so $m_i$ is the
posterior predictive of $x_i$. (For a coin with a uniform prior,
$m_i(\text{heads}) = (k+1)/(i+1)$ after $k$ heads in $i-1$ tosses.) Each
factor has conditional mean $\int m_i = 1$ under the null, whatever the prior;
$\log \mathrm{BF}_t$ is the difference between the cumulative log scores of
the learner and of the null \citep{dawid1984present}; and for i.i.d.\ data
and a regular model with $d$ parameters
\[
\log \mathrm{BF}_t = \max_\theta \sum_{i \le t} \log\frac{q_\theta(x_i)}{p_0(x_i)}
- \frac d2 \log t + O_P(1)
\]
\citep{clarke1990information}. The Monte Carlo e-process has the same
structure, and we now look at it from this point of view: the score, through
its tilted null, plays the role of the forecast $m_i$, the simulations that of
the null density $p_0$, and a family of scores $T_\theta$ that of the model
$q_\theta$.

\subsection{Scores are forecasts}\label{sec:scores}

Fix a time $t$ and condition on $\G_t$, so that the action is known. Write
$p_t$ for the density of the simulator's next outcome under that action, and
$q_t$ for the density of the unknown law $Q_t$ of the real next outcome. The
natural score is the density ratio $q_t/p_t$. Neither density can be
evaluated, but their ratio can be estimated by a classifier separating real
outcomes (label 1) from simulator draws (label 0): if the two classes have
prior weights $r$ and $1-r$, Bayes' theorem gives
\[
\Pr(\text{real} \mid x) = \frac{r\, q_t(x)}{r\, q_t(x) + (1-r)\, p_t(x)},
\qquad
\log\frac{\Pr(\text{real} \mid x)}{\Pr(\text{simulated} \mid x)}
= \log\frac{q_t(x)}{p_t(x)} + \log\frac{r}{1-r},
\]
so the classifier's logit is the log ratio up to an additive constant, which
is zero with balanced class weights \citep{thomas2022lfire} and in any case
cancels in \eqref{eq:mc}.

In practice $T_t$ is only an estimate of the true ratio, learned from past
data, and what matters is the density it induces. Reweighting the simulator by
the score as in \eqref{eq:tilt}, which keeps the simulator's outcomes and gives
each a weight proportional to its score, yields the tilted null
\[
\tilde T_t(x) = \frac{T_t(x)\,p_t(x)}{Z_t},
\qquad
Z_t = \E_{P_t}[T_t] = \int T_t(x)\,p_t(x)\,dx ,
\]
a density on the space of outcomes, not the distribution of the score values.
It is the alternative implied by the bet in the sense of
\citet{shafer2021testing}. As $M \to \infty$, $\widehat Z_t \to Z_t$ by the
law of large numbers, so the bet tends to
\[
E_t = \frac{T_t(X_t)}{\widehat Z_t} \;\to\; \frac{T_t(X_t)}{Z_t}
= \frac{\tilde T_t(X_t)}{p_t(X_t)} ,
\]
a likelihood ratio of $\tilde T_t$ against the simulator: a bettor using the
score acts as if the next outcome followed $\tilde T_t$, and wins when real
outcomes fall where $\tilde T_t$ puts more mass than the simulator. If $T_t$
is the true ratio, $Z_t = 1$ and $\tilde T_t = q_t$; otherwise $\tilde T_t$
is the score's forecast of $q_t$. Conversely, any predictive density $f_t$
computed from the past is the tilted null of a score,
\[
T_t = \frac{f_t}{p_t}
\quad\Longrightarrow\quad
Z_t = 1,
\qquad
\tilde T_t = f_t ,
\]
and multiplying a score by a constant $c > 0$ changes neither the bet
\eqref{eq:mc}, where $c$ cancels, nor its tilted null. Scores up to a
constant and predictive densities are therefore the same object. The
correspondence is the familiar one between density ratios and densities,
used for instance in noise-contrastive estimation \citep{gutmann2012noise};
what the Monte Carlo normalisation adds is that the wealth scores these
densities prequentially, without evaluating any of them.

Indeed, by the limit of the bet above, with the betting fraction
$\lambda = 1$ the log wealth tends to
\begin{equation}\label{eq:prequential-wealth}
\log W_t \;=\; \sum_{i \le t} \log E_i \;\to\; \sum_{i \le t} \log \tilde T_i(X_i)
- \sum_{i \le t} \log p_i(X_i),
\end{equation}
which has the form of $\log \mathrm{BF}_t$ in \eqref{eq:prequential}, with
the tilted null $\tilde T_i$ in place of the posterior predictive $m_i$ and
the simulator $p_i$ in place of the null density $p_0$. It is the difference
between the cumulative log scores of two forecasters, each forecast issued
before the outcome it is scored on: the prequential likelihood ratio of
\citet{dawid1984present}. The monitor rejects at level $\alpha$ when the
score's forecasts have predicted the real outcomes $1/\alpha$ times better
than the simulator did. Under the truth, writing
$\log(\tilde T_t/p_t) = \log(q_t/p_t) - \log(q_t/\tilde T_t)$, the expected
increment of \eqref{eq:prequential-wealth} is
\[
\E_{Q_t}\Big[\log\frac{\tilde T_t(X_t)}{p_t(X_t)}\Big]
= \KL(Q_t\|P_t) - \KL(Q_t\|\tilde T_t).
\]
The first term is the information against the simulator and does not depend
on the score; the second is the error of the score's forecast, which by
Gibbs' inequality is nonnegative and vanishes only when $\tilde T_t = Q_t$,
that is, when $T_t$ is the true ratio up to a constant. This is
Theorem~\ref{thm:growth}(ii), and a finite $M$ multiplies the growth by about
$M/(M+1)$ (Theorem~\ref{thm:growth}(iv)). 

Reading the score as a forecast scored by the logarithmic loss turns the
choice of the score into a problem of sequential prediction, which
Section~\ref{sec:rules} solves with a Bayesian mixture. The reading also says
where power comes from. A score can only forecast within the family of its
tilted nulls, so the features determine which departures can be detected: for
a score built from the simulator's own parameters, as in SLICE, the tilted
nulls belong to that parametric family, and a departure outside it is
invisible however accurate the score (Section~\ref{sec:exp-concurrent}),
whereas a score learned on features of the outcome reweights the simulator in
directions that need not correspond to any parameter. 

\subsection{From Bayes to the simulation likelihood}\label{sec:rules}\label{sec:exactbayes}\label{sec:simlik}

An update rule for the score must fix the bet on
$X_t$ before $X_t$ is seen, should hedge early on, since an overconfident
forecast can lose all the growth (Theorem~\ref{thm:growth}), and must use
simulations alone, since $p_t$ can only be sampled.

Bayes' rule hedges by averaging over what has not yet been learned. If the
real outcomes followed a model $\{Q_\theta\}$ with prior $\pi_0$, the best
forecast under the model would be its posterior predictive $m_t$, since the
logarithmic score is strictly proper. Section~\ref{sec:scores} supplies a
model: the log-linear scores $s = \theta^\top\phi$ have tilted nulls
$dQ_\theta/dP = \exp\{\theta^\top\phi - \log Z(\theta)\}$, with
$Z(\theta) = \E_P[e^{\theta^\top\phi}]$, an exponential family built on the
simulator. With this model the forecast is an average of tilted nulls,
$m_t = \int q_\theta\,\pi_{t-1}(d\theta)$, its score is
$T_t = \int (dQ_\theta/dP)\,\pi_{t-1}(d\theta)$, and the process becomes a
Bayes factor, the benchmark for what follows.

\begin{proposition}[Predictive scores give the Bayes factor; cf.\ \citealp{clarke1990information}]\label{prop:bayes}
With this score $\E_P[T_t] = 1$, the unlimited-budget factor is
$E_t^\infty = T_t(X_t)$, and
\[
\prod_{i\le t} E_i^\infty = \int \prod_{i \le t} \frac{dQ_\theta}{dP}(X_i)\,\pi_0(d\theta),
\]
the Bayes factor of the tilted family against the simulator. For i.i.d.\ data
and under the usual regularity conditions for Laplace approximations, its
logarithm is $\sup_\theta \sum_{i\le t}\log (dQ_\theta/dP)(X_i) - \frac{d}{2}\log t
+ O_P(1)$, and $t^{-1}$ times it converges to
$\KL(Q\|P) - \inf_\theta \KL(Q\|Q_\theta)$.
\end{proposition}

The proof, in Section~S7 of the Supplement, combines the chain rule with the
Laplace expansion of \citet{clarke1990information}.

Proposition~\ref{prop:bayes} is only a benchmark: Bayes' rule needs the
normalising constant $Z(\theta)$ of every tilt at every step, an integral
against a simulator that can only be sampled.

The Monte Carlo factor provides the way out, because it estimates that
constant from the $M+1$ outcomes. Fix $\theta \in \mathbb{R}^d$ and
$\G_t$-measurable features $\phi_t$, write $U_{t,0} = X_t$ and
$U_{t,k} = Y_{t,k}$, and let
\[
E_t(\theta) = \frac{(M+1)\exp\{\theta^\top\phi_t(U_{t,0})\}}
{\sum_{k=0}^{M}\exp\{\theta^\top\phi_t(U_{t,k})\}}
\]
be the factor \eqref{eq:mc} with score $s_t = \theta^\top\phi_t$. As a
function of $\theta$ it is, up to the constant $M+1$, the probability that the
score assigns to the real outcome among the $M+1$
(Section~\ref{sec:method}), a conditional-logit likelihood whose normalising
constant is a sum over the simulations, always computable. We call
$L_t(\theta) = \prod_{i \le t} E_i(\theta)$ the \emph{simulation likelihood};
its logarithm is the InfoNCE objective of contrastive learning
\citep{oord2018representation,ma2018noise}, which \citet{durkan2020contrastive}
use to estimate likelihood ratios in simulation-based inference. With a prior $\pi_0$ on $\theta$
and posterior $\pi_t \propto \pi_0 L_t$, the \emph{posterior-predictive bet}
is
\[
\bar E_t = \int E_t(\theta)\, \pi_{t-1}(d\theta),
\]
updated by one posterior step per observation instead of a refit on the whole
history. The posterior is a prequential posterior in the sense of
\citet{sinharoy2026prequential}, with loss $-\log E_t(\theta)$ and learning
rate one, the choice that makes the posterior mixture of the bets an
e-process.

\begin{theorem}[Predictive updating on the simulation likelihood]\label{thm:simlik}
Suppose the conditions of Theorem~\ref{thm:valid} hold for each fixed $\theta$
(simulations under the action taken, $\phi_t$ $\G_t$-measurable), that
$\|\phi_t\| \le G$, and that $\pi_0$ is the uniform distribution on a convex
compact set $\Theta \subset \mathbb{R}^d$ with nonempty interior and diameter
at most $2D$.
\begin{enumerate}
  \item[(i)] (Exact factorisation and validity.) The wealth
        $\bar W_t = \prod_{i \le t} \bar E_i$ equals
        $\int L_t(\theta)\,\pi_0(d\theta)$. It is a test martingale under
        $H_0$, for every policy, and it is updated by one posterior step per
        observation, with no recomputation.
  \item[(ii)] (Pathwise regret.) On every sample path, for every
        $t \ge d/(4GD)$,
        \[
        \log \bar W_t \;\ge\; \sup_{\theta \in \Theta} \log L_t(\theta)
        - d \log t - d \log\frac{4eGD}{d}.
        \]
        If the maximiser $\theta^*_t$ of $L_t$ lies in $\Theta$ at distance at
        least $r_t = \sqrt{d/(tG^2)}$ from its boundary, the bound improves to
        $\log \bar W_t \ge \log L_t(\theta_t^*) - \frac d2 \log t
        - \frac d2 \log\frac{e G^2 (2D)^2}{d}$.
  \item[(iii)] (Growth.) Let $g_{M,i}(\theta) = \E[\log E_i(\theta) \mid \G_i]$,
        the conditional growth rate \eqref{eq:growthrate} of the fixed score
        $T = e^{\theta^\top\phi}$, and $g_M(\theta) = g_M(Q, P, e^{\theta^\top\phi})$.
        For every $\theta \in \Theta$, almost surely,
        \[
        \liminf_{t \to \infty} \frac1t \Big(\log \bar W_t - \sum_{i \le t} g_{M,i}(\theta)\Big) \ge 0 .
        \]
        In particular, for i.i.d.\ data under a fixed action and features,
        $\liminf_t t^{-1}\log \bar W_t \ge \sup_{\theta\in\Theta} g_M(\theta)$,
        and, if $\KL(Q\|P) < \infty$, by Theorem~\ref{thm:growth}
        $g_M(\theta) \ge \KL(Q\|P) - \KL(Q\|Q_\theta) - \log\{1 + (\rho_\theta - 1)/(M+1)\}$,
        where $Q_\theta$ is the tilt of $P$ by $\theta^\top\phi$ and
        $\rho_\theta = \E_Q[e^{\theta^\top\phi}]/\E_P[e^{\theta^\top\phi}]$.
\end{enumerate}
\end{theorem}

Part (i) says that the score can be updated by Bayes' rule exactly, although
the normalising constant $Z(\theta)$ of the tilt is unknown. The simulation
likelihood is not the density of the real outcome under the tilt; it answers
the question of Section~\ref{sec:method}, which of the $M+1$ outcomes at hand
is the real one, and needs only their scores. Its posterior predictive is the
bet itself, so the guarantees of Bayesian mixtures hold for the wealth, and as
$M \to \infty$ it recovers the Bayes factor of Proposition~\ref{prop:bayes}.

Part (ii) holds on every path, with no assumption on the alternative: the
monitor is never more than $d\log t$ behind the best fixed score in the class
chosen in hindsight and evaluated on the same simulations. The argument is the
standard one for Bayesian mixtures of log-concave likelihoods
\citep{kakade2005online,grunwald2007minimum}, applied to the simulation
likelihood, whose log is concave in $\theta$ and whose gradient is bounded by
$2G$. Part (iii) turns this into growth at the rate of the best tilt in the
class, up to the simulation cost of Theorem~\ref{thm:growth}(iii), and
Ville's inequality then gives rejection after about
$\{\log(1/\alpha) + d\log t\}/\sup_\theta g_M(\theta)$ steps.

\subsection{Practical update rules}\label{sec:blr}\label{sec:algorithm}

Exact Bayesian updating needs the posterior $\pi_{t-1}$ at every step. We
replace it by a Gaussian approximation updated in one step per observation,
on either of two losses, or update the forecast directly as martingale
posteriors do. All give valid bets; they differ in where uncertainty comes from
and in the guarantees they carry (Table~S8 of the Supplement).

The integral in $\bar E_t$ is over $\theta$ only, and any $\G_t$-measurable
distribution $\nu_{t-1}$ on $\theta$, an approximate posterior, in place of
$\pi_{t-1}$ still gives a valid bet
$E^{\nu}_t = \int E_t(\theta)\,\nu_{t-1}(d\theta)$, because $E_t(\theta)$ is
fair for each $\theta$. The price of the approximation is bounded as follows.

\begin{proposition}[Approximate predictive bets]\label{prop:approx}
Let $\nu_{t-1}$ be any $\G_t$-measurable distribution on $\mathbb{R}^d$ with mean
$\bar\theta_{t-1}$ and covariance $\Sigma_{t-1}$, and $\|\phi_t\| \le G$. Then
$E^{\nu}_t$ is an e-value given $\G_t$, and
\[
\log E^{\nu}_t \;\ge\; \E_{\nu_{t-1}}\big[\log E_t(\theta)\big]
\;\ge\; \log E_t(\bar\theta_{t-1}) - \frac{G^2}{2}\,\mathrm{tr}\,\Sigma_{t-1}.
\]
\end{proposition}

\begin{proof}
Validity follows from Fubini's theorem. The first inequality is Jensen's.
The function $\theta \mapsto \log E_t(\theta)$ is concave, with gradient
$\phi_t(U_{t,0}) - \sum_k w_k \phi_t(U_{t,k})$ and Hessian
$-\mathrm{Cov}_w(\phi_t)$, where $w_k \propto e^{\theta^\top\phi_t(U_{t,k})}$;
since $\mathrm{Cov}_w(\phi_t) \preceq \E_w[\phi_t\phi_t^\top] \preceq G^2 I$, with $\E_w$
the average over $k$ of $\phi_t(U_{t,k})$ with weights $w_k$,
Taylor's theorem with integral remainder around $m = \bar\theta_{t-1}$ gives
$\log E_t(\theta) \ge \log E_t(m) + \nabla\log E_t(m)^\top(\theta - m)
- \frac{G^2}{2}\|\theta - m\|^2$ for every $\theta$. Taking the expectation
under $\nu_{t-1}$, the linear term vanishes because $m$ is its mean, and
$\E_{\nu_{t-1}}\|\theta - m\|^2 = \mathrm{tr}\,\Sigma_{t-1}$.
\end{proof}

A predictive bet built from an approximate posterior is therefore never worse
than the plug-in bet at its mean by more than
$\frac{G^2}{2}\sum_i \mathrm{tr}\,\Sigma_{i-1}$. If the posterior precision grows
at least linearly, $\Sigma_i^{-1} \succeq \Sigma_0^{-1} + i\,c_0 I$ for some
$c_0 > 0$, this is at most $\frac{G^2 d}{2c_0}\log t + O(1)$, the order of the
regret in Theorem~\ref{thm:simlik}; such growth needs features that are not
collinear under the simulator and iterates that stay bounded. The Bayesian
learning rule \citep{khan2023bayesian} provides such a posterior. It
approximates a posterior within an exponential family by natural-gradient
descent on the expected loss plus the Kullback--Leibler divergence from the
prior; in a stream, one step per observation with the previous approximation
as the prior \citep{jones2024bong} gives, for a Gaussian
$\nu_t = \mathcal N(\bar\theta_t, S_t^{-1})$ and a per-step loss $\ell_t$,
\begin{equation}\label{eq:blr}
S_t = S_{t-1} + \E_{\nu_{t-1}}[\nabla^2 \ell_t], \qquad
\bar\theta_t = \bar\theta_{t-1} - S_t^{-1}\, \E_{\nu_{t-1}}[\nabla \ell_t],
\end{equation}
which is exact Bayesian updating (the Kalman filter) for a Gaussian likelihood
and one step of online Newton in general. Its cost is $O((M+1)d^2 + d^3)$ per
step with a direct solve, $O((M+1)d^2)$ with rank-one updates of $S_t^{-1}$,
with no refitting and no burn-in.

Two losses are natural. On the simulation likelihood,
$\ell_t(\theta) = -\log E_t(\theta)$, with gradient
$\sum_k w_k \phi_t(U_{t,k}) - \phi_t(U_{t,0})$ and Hessian
$\mathrm{Cov}_w(\phi_t)$; expectations under $\nu_{t-1}$ are computed with
fixed antithetic nodes, and the bet is the average of $E_t(\theta)$ over the
nodes, so Theorem~\ref{thm:simlik} and Proposition~\ref{prop:approx} apply
directly. Alternatively, $\ell_t$ is the balanced logistic loss of classifying
the real outcome (weight one) against the $M$ simulations (weight $1/M$ each),
a generalised posterior \citep{bissiri2016general,sinharoy2026prequential}
that identifies the same tilt through noise-contrastive estimation
\citep{gutmann2012noise}. With the probit approximation
$\E_\nu[\sigma(z^\top\theta)] \approx \sigma(\kappa\bar\eta)$, where
$\bar\eta = z^\top\bar\theta$, $v = z^\top S^{-1}z$ and
$\kappa = (1+\pi v/8)^{-1/2}$, the score is the predictive log-odds
$\kappa\bar\eta$ on the probability scale, inserted in \eqref{eq:mc}; this is
our default. The ratio scale $\bar\eta + v/2$ would be the exact predictive
under an exact posterior, but its premium $v/2$ over-bets where the
approximation is least reliable, and it is the worst score in every experiment
(Supplement, Section~S4). Both losses give scores that start at zero, so the
monitor starts by not betting, and that shrink where the posterior is
uncertain. Features are standardised with simulator draws only, before any
real data are seen, so that the score stays predictable.

The same two classifiers can also be updated without a posterior over their
weights, by updating their predictive directly. In the approach of
\citet{fong2023martingale} one specifies a rule
$m_{t+1} = \mathcal U_t(m_t, X_t)$ and recovers uncertainty by
\emph{predictive resampling}, drawing future outcomes from the current
predictive and updating; this converges when imputing one more outcome does
not move the forecast on average, $\E_{X \sim m_t}[\mathcal U_t(m_t, X)] = m_t$,
the conditionally identically distributed sequences of \citet{berti2004limit}.
Here the learner carries the predictive itself, $\sigma(z^\top\theta_t)$ for
the logistic loss or the softmax $e^{\theta_t^\top z_k}/\sum_{k'} e^{\theta_t^\top z_{k'}}$
over the $M+1$ outcomes for the simulation likelihood, and updates it by
Fisher scoring on the predictive's own log-likelihood,
\begin{equation}\label{eq:mp}
I_t = I_{t-1} + J_t(\theta_{t-1}), \qquad
\theta_t = \theta_{t-1} - I_t^{-1}\, \nabla_t(\theta_{t-1}),
\end{equation}
where $\nabla_t$ is the gradient of the new batch's loss, $J_t$ its Fisher
information, which for both losses does not depend on the labels, and $I_0$ a
prior precision. Under the predictive the gradient has mean zero, so
$\theta_t$ is a martingale, and the bet averages the factor over $N$
predictive-resampling paths, obtained by running \eqref{eq:mp} forward from
$(\theta_{t-1}, I_{t-1})$ on past batches whose labels are drawn from the
predictive ($N = 16$, truncated at $\min(t, 200)$ steps, refreshed every ten
steps). Any learner can be made a martingale-posterior learner in this way:
with the empirical predictive, predictive resampling is the Bayesian
bootstrap, and a random forest fitted with Dirichlet weights on the past steps
is valid as well (Supplement, Section~S4).

Algorithm~S1 in the Supplement puts the default rule together step by step,
with a diagram of one step of the procedure.

%======================================================================
\section{Simulation studies}\label{sec:experiments}
%======================================================================

The experiments check validity with learned scores and adaptive allocation,
the theory of Sections~\ref{sec:power} and~\ref{sec:predictive}, the update
rules and the concurrent methods, and the gain from adaptive actions in an
audit with real model outputs. Throughout, $\alpha = 0.05$,
$M = 16$ unless stated otherwise, and the betting fraction is mixed over 20
values from 0.05 to 1. Each table reports, over 400 replications on data
shared by all methods, the share of runs that reject a true null within the
horizon (false alarms; Monte Carlo standard error about 0.011 near 0.05), the
share that reject a false null (power), and the median stopping time, with
runs that never reject counted as infinite; the horizon is 500 steps for the
drift task and 400 for the bandit and the audit.

\subsection{Validity with learned scores}\label{sec:exp-drift}

The simulator is the two-dimensional ``banana''
$x = (z_1, z_2 + b(z_1^2 - 1))$, $z \sim \mathcal N(0, I)$, with $b = 0.5$; under
the alternative the real data drift to $b = 0.8$, which changes neither the
means nor the marginal of $x_1$, so the shape of the discrepancy must be
learned. The score is a balanced polynomial logistic regression refitted every
ten steps, regularised (degree 4, ridge 1) or overfitting (degree 6, ridge
$10^{-3}$). With the same score, using the learned ratio directly as the bet
gives 10\% false alarms with the regularised learner and 59\% with the
overfitting one, while the Monte Carlo e-process gives 3\% with both
(Table~\ref{tab:drift}): a worse learner costs power, not validity. A
Kolmogorov--Smirnov test recomputed every ten steps gives 32\%. The oracle
likelihood-ratio e-process, unavailable in practice, stops after 35 steps;
the gap to the 198 steps of the learned score is the price of learning the
discrepancy online.

\begin{table}[tp]
\spacingset{1}
\centering
\caption{Generative-model drift (Section~\ref{sec:exp-drift}). The KS test is
applied per coordinate against a pool of simulator draws, with a Bonferroni
correction, and recomputed every ten steps.}
\label{tab:drift}
\begin{tabular}{llrrr}
\toprule
learner & method & false alarms & power & median stop \\
\midrule
regularised & MC e-process & 0.030 & 0.990 & 198 \\
 & plug-in & 0.102 & 0.995 & 158 \\
flexible & MC e-process & 0.028 & 0.838 & 320 \\
 & plug-in & 0.588 & 0.998 & 66 \\
-- & KS, recomputed & 0.323 & 0.750 & 250 \\
-- & oracle LR & 0.030 & 1.000 & 35 \\
\bottomrule
\end{tabular}
\end{table}

\subsection{Adaptive allocation}\label{sec:exp-bandit}

Three arms are allocated by Thompson sampling on the real outcomes, and a twin
simulates any arm, $y = \mu_k + e^{\epsilon/2} - e^{1/8}$ with
$\epsilon \sim \mathcal N(0, 1)$ and $\mu = (0, 0.2, 0.4)$. Under the
alternative arm~0 is better in reality than the twin predicts (mean shifted by
0.5), and Thompson sampling pulls it in 65\% of steps instead of 7\%. The
conditioned e-process stays valid, with 2.8\% false alarms
(Table~\ref{tab:bandit}), whereas the same e-process with the twin run on
uniformly random arms, the setting of Proposition~\ref{prop:necessary}, rejects
the true null in every run, and a per-arm $t$-test recomputed every ten steps
in 39\% of runs.

\begin{table}[tp]
\spacingset{1}
\centering
\caption{Thompson-sampling bandit against a digital twin
(Section~\ref{sec:exp-bandit}).}
\label{tab:bandit}
\begin{tabular}{lrrr}
\toprule
method & false alarms & power & median stop \\
\midrule
MC e-process, twin run on the chosen arm & 0.028 & 1.000 & 49 \\
MC e-process, twin run on random arms & 1.000 & 1.000 & 40 \\
plug-in & 0.055 & 1.000 & 44 \\
per-arm $t$-test, recomputed & 0.388 & 1.000 & 30 \\
oracle LR & 0.043 & 1.000 & 15 \\
\bottomrule
\end{tabular}
\end{table}

\subsection{Empirical growth rates}\label{sec:exp-budget}

On the drift task, validity holds at every budget, including $M = 1$ (false
alarms 1.5--3.5\%), and the median stopping time at $M = 1$ is 1.8 times that
at $M = 64$, close to the factor 2 of Theorem~\ref{thm:growth}(iv). With
$P = \mathcal N(0,1)$, $Q = \mathcal N(\delta, 1)$ and $T(x) = e^{\theta x}$,
where everything is explicit, the bounds of Theorem~\ref{thm:growth} hold in
all 56 configurations, and $g_M/g_\infty$ follows $M/(M+1)$ closely for
perfect, under- and overconfident scores alike. The regret bound of
Theorem~\ref{thm:simlik} holds on every path, with largest regret 7.5 against
a bound of 21.5 at $t = 2000$ (Supplement, Section~S4, with the figure).

\subsection{Update rules}\label{sec:exp-blr}

All the update rules of Section~\ref{sec:predictive}, updated at every step
with no burn-in, give valid tests, with 1.3 to 3\% false alarms, whereas using
the logistic posterior odds directly as the bet gives 14 to 74\%. The logistic
Bayesian learning rule is the fastest and makes the method insensitive to the
flexibility of the features (power 0.97 with degree-6 features, against 0.84
for the refit). The simulation-likelihood learner, the one with the guarantee
of Theorem~\ref{thm:simlik}, is about 6\% slower on the drift task: the
Bayesian mixture pays for its pathwise bound with cautious early bets. All
variants are in the Supplement (Section~S4), and a comparison with the
Besag--Clifford e-process of \citet{dombowsky2026besag} in Section~S3.

\subsection{Comparison with concurrent methods}\label{sec:exp-concurrent}

The concurrent methods use the same bet and differ in how the score is
chosen, so we run them with the same simulations, mixture over $\lambda$ and
data. For \citet{forre2026anytime} we use the score refitted every ten steps by
maximising the penalised empirical growth, and a score fitted in advance on
simulations from a guessed alternative. For \citet{cortinovis2026slice} we use
SLICE with a single null: an amortised neural ratio network trained on
$2 \times 10^5$ simulations from the simulator's parametrised family
(curvature $b$ for the drift, a location shift per arm for the bandit),
evaluated at a predictable estimate of the parameter. In the bandit both are
run on simulations under the chosen arm, and once as published, with the twin
on random arms. Each task also has an alternative outside SLICE's family: in
the drift the first coordinate is 20\% wider, in the bandit the most pulled
arm keeps its mean and has 50\% larger noise.

\begin{table}[tp]
\spacingset{1}
\centering
\caption{Comparison with \citet{forre2026anytime} (F\&B) and
\citet{cortinovis2026slice} (SLICE). Entries: power and, in brackets, median
stopping time; ``in family'': the alternatives of
Sections~\ref{sec:exp-drift} and~\ref{sec:exp-bandit}; ``out of family'':
departures that SLICE's family does not contain. False alarms: range over the
two tasks. Time: computing time per replication on one core, drift / bandit,
excluding the training of SLICE's network. All variants are in the
Supplement.}
\label{tab:concurrent}
\small
\begin{adjustbox}{max width=\textwidth}
\begin{tabular}{lrrrrrr}
\toprule
 & \multicolumn{2}{c}{drift} & \multicolumn{2}{c}{bandit} & false & time \\
score & in family & out of family & in family & out of family & alarms & (s) \\
\midrule
ours, BLR logistic & 0.99 (214.5) & 0.73 (321.5) & 1.00 (38) & 1.00 (112) & 2.0--2.2\% & 0.3 / 0.2 \\
ours, simulation likelihood & 0.96 (228.5) & 0.87 (228) & 1.00 (55) & 0.99 (157) & 1.0--1.8\% & 0.3 / 0.15 \\
F\&B, refitted on the past & 0.99 (210) & 0.91 (228) & 1.00 (61) & 0.99 (167) & 1.5--1.8\% & 1.0 / 0.65 \\
F\&B, fitted in advance & 1.00 (57.5) & 0.01 & -- & -- & 4.8\% & -- \\
SLICE, neural ratio, $G = 1$ & 1.00 (66.5) & 0.01 & 1.00 (35) & 0.56 (258) & 1.5--3.3\% & 0.45 / 0.4 \\
F\&B, twin on random arms & -- & -- & \multicolumn{2}{c}{false alarms 100\%} & -- & -- \\
\bottomrule
\end{tabular}
\end{adjustbox}
\end{table}

In Table~\ref{tab:concurrent}, conditioning on the action is necessary for
every score: run as published, the refitted score of
\citet{forre2026anytime} rejects the true null in every bandit run. Scores
that know where the departure lies are fastest inside the family, SLICE after
67 steps on the drift and a score fitted on the true alternative after 58,
against about 210 for scores learned from the data; but outside the family
their power falls to 0.01 on the drift (Remark~\ref{rem:family}), while every
score learned from real against simulated outcomes keeps it. Among the learned
scores, the refit of \citet{forre2026anytime} is slightly more powerful on the
drift task, markedly so against our logistic learner outside the family (0.91
against 0.73), and ours are faster on the bandit (38 against 61 steps inside
the family, 112 against 167 outside). The gap outside the family comes from
the logistic learner, not from the method. On the simulation likelihood, the
objective that the refit maximises, our update has the same median stopping
time as the refit (228 steps), plausibly because the logistic learner's probit
shrinkage keeps its bets small for longer when the signal is weak. Ours cost a
third as much and carry the pathwise guarantee of Theorem~\ref{thm:simlik}.
They average over the uncertainty about the score's parameter, whereas a sequence of refits at a maximiser is in general not the
predictive sequence of any exchangeable model
\citep{fortini2000exchangeability,fortini2024exchangeability}.

\begin{remark}[Inside and outside the family]\label{rem:family}
SLICE bets with $s_t = \log (p_{\hat\theta_{t-1}}/p_{\theta_0})$ within a
parametrised family $\{P_\theta\}$ with null $\theta_0$, so its tilted null is
$P_{\hat\theta_{t-1}}$. Outside the family, $\hat\theta_t$ converges to
$\theta^* = \arg\min_\theta \KL(Q\|P_\theta)$, and by
Theorem~\ref{thm:growth}(ii) the growth tends to
$\KL(Q\|P_{\theta_0}) - \KL(Q\|P_{\theta^*})$, which is zero when
$\theta^* = \theta_0$, however large the departure. Both out-of-family
alternatives are of this kind, whereas our features contain $x_1^2$ and the
squared outcome of each arm, so a change of scale moves the tilted null towards
$Q$. Details, and the choice between the NRE-A and NRE-C estimators
\citep{miller2022contrastive}, are in the Supplement.
\end{remark}

\subsection{Auditing a language-model API with adaptively chosen prompts}\label{sec:exp-llm}

Many language models are used through an API, a service that returns the
model's reply to a text, the prompt, without showing which model produced it.
A provider may advertise a model $P$ and answer some queries with a cheaper
substitute, such as $P$ at lower numerical precision or a smaller model. An
audit checks this from the outside, comparing the API's replies with replies
of $P$ to the same prompts, run locally. The API is the real system, the local
copy of $P$ is the simulator and the prompt is the action; neither gives
probabilities, only replies.

As in the routing audit of \citet{cortinovis2026slice}, the prompts ask for
random output (eight random digits, an integer between 1 and 100, a colour,
six coin flips), so that a substitute shows up as a change in the frequencies
of the replies rather than as a wrong answer. The advertised model is
Qwen2.5-1.5B-Instruct \citep{qwen2024qwen25} stored with 8 bits per weight;
the substitutes are the same model with 4 and 3 bits and the smaller
Qwen2.5-0.5B-Instruct. All replies are real outputs of \texttt{llama.cpp}
\citep{gerganov2023llamacpp} at temperature one; only the routing, the choice
of the queries that go to the substitute, is simulated. Either 30\% of all
queries go to the substitute, or 50\% of the colour and number queries and
none of the others. In a last scenario the API answers every query with a
sharpened $P$, which gives $P$'s frequent replies more often, as a lower
temperature would: the most frequent colour goes from 23\% to 44\% of the
replies.

At each step the auditor chooses a prompt, sends it to the API and generates
$M = 8$ replies of $P$ to the same prompt. The prompt is chosen adaptively,
from the bets won so far. For each prompt the auditor computes the average log
growth of its past bets, adds a bonus that shrinks as the prompt is asked more
often, and asks the prompt with the largest sum, so that it tries prompts it
knows little about as well as those that have paid; this is the upper
confidence bound (UCB) rule of bandit problems \citep{lattimore2020bandit},
used here without tuning. No prompt is asked in fewer than a tenth of the
queries: whenever one falls below that share, it is asked next, so that the
score keeps learning on every prompt.

We compare five scores. Two are ours, both on hashed character $n$-grams
with one block of weights per prompt: the logistic BLR of
Section~\ref{sec:blr}, our default, and the simulation-likelihood rule of the
same section, the one covered by Theorem~\ref{thm:simlik}. They reweight the
replies of $P$ prompt by prompt and need no reply of the substitute. The score
of \citet{forre2026anytime} is used in its two versions: refitted every ten
steps on the whole history by maximising the empirical growth, on the same
features, and fitted in advance to separate a guessed substitute $R$ from
$P$. SLICE assumes that $R$ answers a fraction $\theta$ of the queries, the
same for all prompts, and estimates $\theta$. The score fitted in advance and
SLICE learn $R$ before the audit from 500 replies per prompt of the 4-bit
model, so they guess wrong for the 3-bit and 0.5B substitutes; routing that
depends on the prompt lies outside SLICE's family, and the sharpened $P$ lies
outside it in a stronger sense, since it is closer to $P$ than to $R$.

\begin{table}[tp]
\spacingset{1}
\centering
\caption{Auditing a language-model API (Section~\ref{sec:exp-llm}): median
number of queries to an alarm (fewest in bold), rounded to an integer, over
the runs that reject within 400 queries, with the power when below one;
``never'': no run rejects. Routing: ``all'', the substitute answers 30\% of
all queries (the sharpened $P$ answers every query); ``two'', it answers 50\%
of the colour and number queries only. SLICE and the score fitted in advance
are given replies of the 4-bit substitute; the refitted score and ours are
given none. False alarms: rejections of the true null within 400 queries
(Monte Carlo standard error about 1.1\%); a twin that answers a random prompt
raises a false alarm in every run. $^\ddagger$Average of the wealths of ours
(logistic) and of the score fitted in advance. $^*$For the exact posterior;
the one-step Gaussian update used here is covered by
Proposition~\ref{prop:approx}.}
\label{tab:llm}
\small
\setlength{\tabcolsep}{4pt}
\begin{adjustbox}{max width=\textwidth}
\begin{tabular}{llrrrrrl}
\toprule
 & & \multicolumn{4}{c}{substitute} & false & growth \\
\cmidrule(lr){3-6}
score & routing & 4-bit & 3-bit & 0.5B & sharpened $P$ & alarms & guarantee \\
\midrule
\multicolumn{8}{l}{\emph{prompts drawn uniformly (as published)}} \\
SLICE & all & 82 & 149 (0.96) & 110 (0.99) & never & 0.0\% & -- \\
 & two & 76 &  & 107 &  & & \\
\addlinespace[3pt]
F\&B, fitted in advance & all & 44 & 72 & 70 & never & 3.5\% & -- \\
 & two & 39 &  & 59 &  & & \\
\addlinespace[3pt]
F\&B, refitted & all & 251 (0.83) & 170 (0.98) & 102 & 93 & 3.8\% & -- \\
 & two & 200 (0.93) &  & 75 &  & & \\
\addlinespace[3pt]
ours, logistic BLR & all & 255 (0.72) & 209 (0.96) & 123 & 86 & 5.3\% & -- \\
 & two & 219 (0.83) &  & 77 &  & & \\
\addlinespace[3pt]
ours, simulation likelihood & all & 254 (0.69) & 195 (0.99) & 100 & 151 & 4.3\% & Thm~\ref{thm:simlik}$^*$ \\
 & two & 219 (0.85) &  & 62 &  & & \\
\midrule
\multicolumn{8}{l}{\emph{prompts chosen by UCB (our design)}} \\
SLICE & all & 72 & 109 (0.97) & 95 & never & 0.0\% & -- \\
 & two & 57 &  & 57 &  & & \\
\addlinespace[3pt]
F\&B, fitted in advance & all & \textbf{41} & \textbf{50} & \textbf{48} & never & 2.5\% & -- \\
 & two & \textbf{30} &  & \textbf{23} &  & & \\
\addlinespace[3pt]
F\&B, refitted & all & 177 (0.98) & 131 (0.98) & 94 & \textbf{78} & 3.3\% & -- \\
 & two & 107 &  & 51 &  & & \\
\addlinespace[3pt]
average$^\ddagger$ & all & 50 & 66 & 56 & 95 & 2.8\% & -- \\
 & two & 39 &  & 29 &  & & \\
\addlinespace[3pt]
ours, logistic BLR & all & 218 (0.96) & 162 (0.99) & 109 & \textbf{78} & 3.0\% & -- \\
 & two & 136 &  & 52 &  & & \\
\addlinespace[3pt]
ours, simulation likelihood & all & 193 (0.92) & 138 (0.99) & 83 & 130 & 1.5\% & Thm~\ref{thm:simlik}$^*$ \\
 & two & 119 (0.99) &  & 36 &  & & \\
\bottomrule
\end{tabular}
\end{adjustbox}
\end{table}

Choosing the prompts helps every score (Table~\ref{tab:llm}): with our
adaptive design every score needs fewer queries than with uniformly drawn
prompts, in every scenario in which it rejects and up to 61\% fewer, and our
scores gain up to 0.24 in power; the auditor learns to ask for a colour more
often, because that is where the substitutes differ most. Validity requires
the conditioning: comparing the API's reply with replies of $P$ to a random
prompt, rather than to the same prompt, raises a false alarm in every run,
whereas the conditioned monitors raise false alarms in at most 5.3\% of the
runs, within Monte Carlo error of the nominal 5\%.


With our design the score fitted in advance is the fastest in the five
routing scenarios, after 23 to 50 queries against 36 to 218 for the scores
learned during the audit. When the API sharpens $P$, however, a score fitted
to $R$ has a tilted null close to $R$, hence a growth rate close to
$\KL(Q\|P) - \KL(Q\|R) < 0$ by Theorem~\ref{thm:growth}: it never rejects,
nor does SLICE, which estimates a routing fraction of zero and stops betting.
The three scores learned during the audit reject in every run, after 78 to
130 queries. The average of the wealths of the logistic score and of the
score fitted in advance keeps the best of both: it is within 16 queries of
the faster component in the routing scenarios, and rejects after 95 queries
when the guess is wrong.

Our simulation-likelihood rule and the refit of \citet{forre2026anytime} both
learn the score during the audit and need no guess about the substitute. With
our design the refit needs 5 to 10\% fewer queries on the 4-bit and 3-bit
substitutes and 40\% fewer on the sharpened $P$ (78 against 130), whereas our
rule needs 12 and 29\% fewer on the 0.5B substitute (83 against 94, and 36
against 51); our logistic rule ties with the refit on the sharpened $P$. In
exchange our rule carries the guarantee of Theorem~\ref{thm:simlik}, through
Proposition~\ref{prop:approx} for the one-step update, and its computing time per query does not grow with the history: 2.1 and
2.2~ms over 400 and 1600 queries, against 8.0 and 53~ms for the refit, which
revisits the whole history at every refit.

Almost all the running time of the audit goes into the simulator: generating
the $M = 8$ replies of $P$ takes 2 to 8 seconds per query on four CPU cores,
whereas the bet and the update take at most 8~ms. The cost of an alarm is
therefore counted in generated replies: SLICE and the score fitted in advance
prepare 6000 replies before the audit and reach the alarm after 6200 to 6900,
whereas the learned scores prepare 800, only to standardise the features, and
reach it after 1100 to 2500 (Supplement, Section~S4, which also reports the
computing time to the alarm).

%======================================================================
\section{A multi-arm multi-stage trial designed in silico}\label{sec:trial}
%======================================================================

\subsection{The trial and its twin}\label{sec:application}

Insulin-dosing algorithms for type~1 diabetes are designed and tested on the
UVA/Padova simulator \citep{dallaman2007meal,kovatchev2009insilico,dallaman2014uva},
accepted by the US Food and Drug Administration for preclinical testing. We use
its open-source implementation \texttt{simglucose} \citep{xie2018simglucose},
a 13-dimensional model with randomness in insulin sensitivity, carbohydrate
absorption and the CGM error \citep{breton2008analysis}; the state is hidden,
so the law of a day's CGM trace has no tractable density.

The trial compares a control strategy (A) with four variants of its bolus,
correction target and suspension threshold (B to E; details in
Section~S8 of the Supplement); the outcome is the time in range (70--180~mg/dL)
over a day with three meals. By the simulator, B and C are worse than control
and D and E equivalent to it. In each of three stages, 50 new participants are
allocated among the active arms, and every enrolled participant is observed for
one more day, carrying their physiological state over. After stages 1 and 2 an
experimental arm is dropped for futility and its participants are re-assigned
to the best active arm. The trial has 150 participants and 300
participant-days, and every element of the design is chosen from the data.

The twin re-initialises the simulator at each participant's current state and
runs it $M = 16$ times under the arm assigned for that day; the score is the
logistic BLR of Section~\ref{sec:blr} on six summaries of the day. We compare
it with the plug-in bet, a \emph{stateless} twin started from rest under the
assigned arm, a \emph{pooled} twin started from rest under a random arm, and a
$t$-test recomputed after every participant-day. To say under which arm the
twin fails, we also run one e-process per arm, betting only on that arm's
participant-days; each is a test martingale by Theorem~\ref{thm:valid}, and
e-BH \citep{wang2022false} across the arms controls the false discovery rate
under the dependence created by the shared interim decisions
\citep{xu2021unified}. Under the null the participants are the simulator. Under
\emph{insulin resistance}, insulin sensitivity is 25\% lower in every arm;
under \emph{defensive snacking}, participants on the intensified arms B and D
eat an unannounced 20~g snack at bedtime.

\subsection{Results}

\begin{table}[tp]
\spacingset{1}
\centering
\caption{A three-stage MAMS trial checked against its in-silico twin (150
participants, 300 participant-days, 100 replications per scenario). Top:
rejection rate and, in brackets, median participant-days to the first alarm.
Middle: share of trials with each arm active at the end. Bottom: share of
trials in which e-BH flags each arm ($\alpha = 0.05$).}
\label{tab:trial}
\begin{adjustbox}{max width=\textwidth}
\begin{tabular}{lrrr}
\toprule
monitor & null & insulin resistance & defensive snacking \\
\midrule
MC e-process, twin at the participant's state & 0.02 & 1.00 (4) & 0.99 (32) \\
plug-in bet & 0.47 & 1.00 (5) & 1.00 (28) \\
stateless twin (ignores the history) & 0.95 & 1.00 (4) & 1.00 (33) \\
pooled twin (ignores allocation and history) & 1.00 & 1.00 (8) & 1.00 (58) \\
$t$-test, recomputed & 0.46 & 1.00 (17) & 1.00 (22) \\
\midrule
\multicolumn{4}{l}{\emph{design: share of trials with the arm active at the end}} \\
B / C / D / E & 0.00 / 0.01 / 0.27 / 0.50 & 1.00 / 0.00 / 1.00 / 0.79 & 0.02 / 0.01 / 0.93 / 0.42 \\
\midrule
\multicolumn{4}{l}{\emph{per-arm e-processes and e-BH: share of trials in which each arm is flagged}} \\
A / B / C / D / E & 0 / 0 / 0 / 0 / 0.01 & 1 / 1 / 1 / 1 / 1 & 0 / 0.64 / 0 / 0.99 / 0.01 \\
\bottomrule
\end{tabular}
\end{adjustbox}
\end{table}

Only the monitor that conditions on both the allocation and the history is
valid, with 2\% false alarms (Table~\ref{tab:trial}). The stateless twin raises
a false alarm in 95\% of trials, because participants no longer start from rest
and re-assigned participants carry the state of a treatment they no longer
receive; the pooled twin does so in every trial, the recomputed $t$-test and
the plug-in bet in 46\% and 47\%. Insulin resistance is detected after a median
of four participant-days, defensive snacking in 99\% of trials after 32. The
snack also changes the design: it protects arm D from night-time lows, so D
survives to the end in 93\% of trials instead of 27\%, for a reason its
in-silico justification does not contain. e-BH flags D in 99\% of trials and B
in 64\%, with a false discovery rate of 0.3\%; B is found less often because
the design usually drops it after about ten participant-days. With the
simulation-likelihood learner the conclusions are the same and detection
slower (51 participant-days instead of 32; Supplement, Section~S4).

The twin is re-initialised at the participant's true state, which in a real
trial would have to be estimated from CGM and insulin records
(Section~\ref{sec:discussion}). An alarm says that the in-silico justification
of the design has failed, not that an arm is unsafe. A single-stage version
with Thompson sampling and the monitoring of a single patient are in the
Supplement.

%======================================================================
\section{Discussion}\label{sec:discussion}\label{sec:related}
%======================================================================

The Monte Carlo e-process separates validity, the control of false alarms,
from power. Validity comes only from how the simulator is run, through
exchangeability with fresh simulations under the action taken and from the
current state, and holds for any score and any policy, including interim
analyses and re-assignments. Power comes only from the score, through the
distance between its tilted null and the truth (Theorem~\ref{thm:growth});
approximations in the score therefore cost power but never validity.

The tilted null also makes a rejection informative. The monitor rejects when
its forecasts, the tilted nulls of its scores, have predicted reality better
than the simulator; as the posterior concentrates, the tilted null approaches
the reweighting of the simulator closest to reality,
$\arg\min_\theta \KL(Q\|Q_\theta)$, and its log ratio to the simulator,
$\theta^\top\phi - \log Z(\theta)$, shows where reality puts more or less mass
than the simulator: a shift in the mean, an excess of variance, a particular
arm. This reading is exploratory rather than inferential. The e-process tests
the global null that the simulator is right, and correlated features can
forecast the same departure equally well, so a large coefficient suggests, but
does not prove, that the simulator fails on that feature. A claim with error
control about one aspect needs its own e-process: a score that depends on the
outcome only through a statistic $f$ is valid for the narrower null that the
simulator reproduces the law of $f$, and several such e-processes, or one per
arm as in the trial, can be combined by e-BH. For a single patient, for instance, monitors on fasting windows and on windows
after declared meals test the twin's basal physiology and its model of meal
absorption separately; a monitor that does not reject gives no evidence that
its part of the twin is right.
The tilted null might also drive the design. Under each action $a$, its
divergence from the simulator, $\KL(\tilde T_{t,a}\|P_{t,a})$, is the growth
the monitor expects from that action if its forecast is right, so choosing the
action that maximises it sends the next observation where the monitor believes
the simulator is most wrong. This uses only the past, so it is valid by
Theorem~\ref{thm:valid}; it is a model-based alternative to the upper
confidence bound of Section~\ref{sec:exp-llm}, which uses the growth actually
realised, and it pays only once the forecast is good.

Monte Carlo tests compare a fixed dataset with simulations from the null
\citep{besag1989generalized}, and \citet{fischer2024sequential} make them
anytime-valid in the number of simulations; here the data themselves arrive
over time, with fresh simulations for each observation. The closest work is
\citet{dombowsky2026besag}, whose results on the likelihood-ratio score are
used, with attribution, in Section~\ref{sec:power}; \citet{forre2026anytime}
show that the likelihood ratio is the optimal score for every $M$, and
\citet{cortinovis2026slice} invert the e-process into confidence sequences for
the parameters of a simulator. All three treat independent data. Relative to
them, ours treats actions chosen from the data and simulators with state,
shows that the conditioning is necessary and quantifies the failure without
it, reads arbitrary scores through the tilted null with the local efficiency
$M/(M+1)$, and updates the score by Bayes' rule on the simulation likelihood,
with a pathwise regret bound, where the concurrent papers bet at a point
estimate.

SLICE \citep{cortinovis2026slice} answers a different question: which values
of a simulator's parameter are compatible with the data. When the departure lies inside its parametrised family it is
the fastest method we compared; when it does not, its power depends on where
the departure projects onto the family and vanishes when the projection is
the null itself (Remark~\ref{rem:family}), and it needs an amortised training
on the family, of the order of $10^7$ simulations in
\citet{cortinovis2026slice}. A score learned from real against simulated
outcomes needs neither, and pays for this with a slower start when prior
knowledge of the alternative exists: in the audit, a score fitted in advance
on the right substitute stops after 30 to 41 queries, the learned score after
136 to 218. The natural synthesis is to use a pretrained score as the centre
of the prior of the online update, so that the monitor starts as fast as the
pretrained score and corrects it where the guess is wrong, with the guarantees
of Theorems~\ref{thm:valid} and~\ref{thm:simlik} relative to that prior; the
average of the two wealths in Section~\ref{sec:exp-llm}, within 16 queries of
the faster component, is a cruder version of it. In the other direction,
since SLICE bets at the plug-in forecast $P_{\hat\theta_{t-1}}$, its confidence
sequences are a likelihood-free version of the running-estimate confidence
sequences of universal inference \citep{wasserman2020universal}. Averaging the
bets over a posterior instead keeps validity, and a pathwise regret bound for
such confidence sequences would need a posterior computed for each candidate
null from its own simulation likelihood, with scores log-linear in their
parameter; we leave this to future work.

Once the simulator is run under the action taken, the design may be chosen to gather
evidence where it accumulates fastest: in the audit this makes every score
faster, ours and the concurrent ones alike, by 8\% to 61\% in the number of
queries, while without the conditioning on the action the concurrent methods
lose validity altogether. Other anytime-valid methods for adaptively collected data, and sequential tests that bet with a learned predictor, are discussed in Section~S10 of the Supplement; none of them uses a simulator as the null. The product of one-step predictives behind
Section~\ref{sec:predictive} has classical forms in prequential analysis,
minimum description length and sequential Monte Carlo (Supplement,
Section~S5); prequential posteriors \citep{sinharoy2026prequential} use
simulations to calibrate a simulator, where we use them to test it.

Several limitations remain. The null is a simple conditional null; simulators
with unknown calibration parameters need composite-null constructions, such
as universal inference or reverse information projections
\citep{grunwald2024safe}, adapted to samplers. Exact validity requires the simulator to be run from the state of the real
system, which is known in silico but must be inferred for a physical twin; an
approximate filter makes exchangeability only approximate. A persistent error
of divergence $\delta$ per step can then raise a false alarm, but not much
before $\log(1/\alpha)/\delta$ steps (Corollary~\ref{cor:falsealarm}(ii)).
Since the resulting bias recurs at every step, extending the guarantees of \citet{barber2026montecarlo} for a
single approximately exchangeable test to the sequential setting is an open
problem. Theorem~\ref{thm:simlik} compares the monitor with the best score in
a parametric class, not with the true density ratio, and its $d\log t$ regret
is not known to be optimal. It holds for the exact posterior, whereas the one-step Gaussian update used in
practice is covered only through Proposition~\ref{prop:approx}; and the
logistic learner, which has no regret guarantee, is faster than the
simulation-likelihood learner, by about 6\% on the drift task and 60\% in the
trial, because the latter bets cautiously. A tempered posterior
$\pi_t \propto \pi_0 L_t^{\gamma}$ with $\gamma > 1$
\citep{sinharoy2026prequential} would bet more and stays valid; a regret bound
for it is the most useful next step.
Theorem~\ref{thm:simlik} and Proposition~\ref{prop:approx} also assume bounded
features, whereas our experiments standardise polynomial features without
bounding them and clip only the score; the guarantees apply exactly once the
features are clipped. Finally, the experiments are low-dimensional;
high-dimensional simulators call for learned feature maps and structured
posterior approximations. High-energy physics is a natural target, since
searches for new physics already compare collision data with large samples
from the Standard Model simulator using learned classifiers
\citep{dagnolo2019learning}, and the Monte Carlo e-process would make such
comparisons anytime-valid as data accumulate.

\bigskip
\begin{center}
{\large\bf SUPPLEMENTARY MATERIAL}
\end{center}

\begin{description}
\item[Supplement:] the notation, the proofs, experimental details, the
timing relative to concurrent work and further related work
(Sections~S6--S10); three
further analyses, of a single patient under closed-loop control, a
single-stage response-adaptive trial, and a comparison with the
Besag--Clifford e-process (Sections~S1--S3); additional results for
Sections~\ref{sec:experiments} and~\ref{sec:trial} (Section~S4); and the
classical forms of the predictive factorisation (Section~S5) (PDF).
\item[Code and data:] Python code for all methods, experiments, tables and
figures, with the replies of the language models of
Section~\ref{sec:exp-llm}\if0\blind, at \url{https://github.com/elenabortolato/evalues4generative}\fi\if1\blind{} (link removed for blind review)\fi. The file \texttt{REPRODUCIBILITY.md} lists,
for every table and figure, the script and command that produce it, the
output file and an indication of the running time. Every replication is
seeded, so the reported numbers are reproduced exactly with the stated library
versions, up to the rounding differences across processors documented there.
\end{description}

\section*{Disclosure statement}
The author reports there are no competing interests to declare.

\section*{Data availability statement}
No external data are used. The replies of the language models in
Section~\ref{sec:exp-llm} were generated from the publicly released Qwen2.5
models and are distributed with the code; all other data are simulated by the
code.

\bibliographystyle{apalike}
\bibliography{references}

\end{document}
